\documentclass[11pt]{article}
% Préambule commun aux cours de la formation C++.
% Inclus par chaque cours.tex via % Préambule commun aux cours de la formation C++.
% Inclus par chaque cours.tex via % Préambule commun aux cours de la formation C++.
% Inclus par chaque cours.tex via \input{preambule}.

\usepackage[utf8]{inputenc}
\usepackage[T1]{fontenc}
\usepackage[french]{babel}
\usepackage{lmodern}
\usepackage{microtype}

\usepackage{geometry}
\geometry{a4paper, margin=2.5cm}

\usepackage{xcolor}
\usepackage{amsmath, amssymb}
\usepackage{enumitem}
\usepackage{hyperref}
\hypersetup{
    colorlinks=true,
    linkcolor=blue!60!black,
    urlcolor=blue!60!black,
    citecolor=blue!60!black,
}

% --- Coloration syntaxique du C++ via listings -----------------------------
\usepackage{listings}

\definecolor{codebg}{rgb}{0.97,0.97,0.95}
\definecolor{codekw}{rgb}{0.10,0.10,0.60}
\definecolor{codecomment}{rgb}{0.20,0.50,0.20}
\definecolor{codestring}{rgb}{0.65,0.15,0.15}
\definecolor{coderule}{rgb}{0.80,0.80,0.80}

% Permet les caractères accentués (UTF-8) à l'intérieur des blocs de code.
\lstset{
    literate=%
        {é}{{\'e}}1 {è}{{\`e}}1 {ê}{{\^e}}1 {ë}{{\"e}}1
        {à}{{\`a}}1 {â}{{\^a}}1 {ä}{{\"a}}1
        {î}{{\^i}}1 {ï}{{\"i}}1
        {ô}{{\^o}}1 {ö}{{\"o}}1
        {ù}{{\`u}}1 {û}{{\^u}}1 {ü}{{\"u}}1
        {ç}{{\c c}}1 {É}{{\'E}}1 {È}{{\`E}}1 {À}{{\`A}}1
}

\lstdefinestyle{cpp}{
    language=C++,
    backgroundcolor=\color{codebg},
    basicstyle=\ttfamily\small,
    keywordstyle=\color{codekw}\bfseries,
    commentstyle=\color{codecomment}\itshape,
    stringstyle=\color{codestring},
    showstringspaces=false,
    numbers=left,
    numberstyle=\tiny\color{gray},
    numbersep=8pt,
    frame=single,
    rulecolor=\color{coderule},
    breaklines=true,
    tabsize=4,
    morekeywords={constexpr,noexcept,override,final,nullptr,auto,
        std,string,vector,size_t,uint64_t,int64_t,optional,namespace}
}
\lstset{style=cpp}

% Raccourci pour du code en ligne, en police machine à écrire.
\newcommand{\code}[1]{\texttt{#1}}

% Boîte « À retenir »
\usepackage{tcolorbox}
\newtcolorbox{aretenir}[1][]{
    colback=yellow!8, colframe=orange!70!black,
    title=#1, fonttitle=\bfseries, sharp corners, boxrule=0.5pt
}
\newtcolorbox{piege}[1][Piège]{
    colback=red!5, colframe=red!60!black,
    title=#1, fonttitle=\bfseries, sharp corners, boxrule=0.5pt
}
.

\usepackage[utf8]{inputenc}
\usepackage[T1]{fontenc}
\usepackage[french]{babel}
\usepackage{lmodern}
\usepackage{microtype}

\usepackage{geometry}
\geometry{a4paper, margin=2.5cm}

\usepackage{xcolor}
\usepackage{amsmath, amssymb}
\usepackage{enumitem}
\usepackage{hyperref}
\hypersetup{
    colorlinks=true,
    linkcolor=blue!60!black,
    urlcolor=blue!60!black,
    citecolor=blue!60!black,
}

% --- Coloration syntaxique du C++ via listings -----------------------------
\usepackage{listings}

\definecolor{codebg}{rgb}{0.97,0.97,0.95}
\definecolor{codekw}{rgb}{0.10,0.10,0.60}
\definecolor{codecomment}{rgb}{0.20,0.50,0.20}
\definecolor{codestring}{rgb}{0.65,0.15,0.15}
\definecolor{coderule}{rgb}{0.80,0.80,0.80}

% Permet les caractères accentués (UTF-8) à l'intérieur des blocs de code.
\lstset{
    literate=%
        {é}{{\'e}}1 {è}{{\`e}}1 {ê}{{\^e}}1 {ë}{{\"e}}1
        {à}{{\`a}}1 {â}{{\^a}}1 {ä}{{\"a}}1
        {î}{{\^i}}1 {ï}{{\"i}}1
        {ô}{{\^o}}1 {ö}{{\"o}}1
        {ù}{{\`u}}1 {û}{{\^u}}1 {ü}{{\"u}}1
        {ç}{{\c c}}1 {É}{{\'E}}1 {È}{{\`E}}1 {À}{{\`A}}1
}

\lstdefinestyle{cpp}{
    language=C++,
    backgroundcolor=\color{codebg},
    basicstyle=\ttfamily\small,
    keywordstyle=\color{codekw}\bfseries,
    commentstyle=\color{codecomment}\itshape,
    stringstyle=\color{codestring},
    showstringspaces=false,
    numbers=left,
    numberstyle=\tiny\color{gray},
    numbersep=8pt,
    frame=single,
    rulecolor=\color{coderule},
    breaklines=true,
    tabsize=4,
    morekeywords={constexpr,noexcept,override,final,nullptr,auto,
        std,string,vector,size_t,uint64_t,int64_t,optional,namespace}
}
\lstset{style=cpp}

% Raccourci pour du code en ligne, en police machine à écrire.
\newcommand{\code}[1]{\texttt{#1}}

% Boîte « À retenir »
\usepackage{tcolorbox}
\newtcolorbox{aretenir}[1][]{
    colback=yellow!8, colframe=orange!70!black,
    title=#1, fonttitle=\bfseries, sharp corners, boxrule=0.5pt
}
\newtcolorbox{piege}[1][Piège]{
    colback=red!5, colframe=red!60!black,
    title=#1, fonttitle=\bfseries, sharp corners, boxrule=0.5pt
}
.

\usepackage[utf8]{inputenc}
\usepackage[T1]{fontenc}
\usepackage[french]{babel}
\usepackage{lmodern}
\usepackage{microtype}

\usepackage{geometry}
\geometry{a4paper, margin=2.5cm}

\usepackage{xcolor}
\usepackage{amsmath, amssymb}
\usepackage{enumitem}
\usepackage{hyperref}
\hypersetup{
    colorlinks=true,
    linkcolor=blue!60!black,
    urlcolor=blue!60!black,
    citecolor=blue!60!black,
}

% --- Coloration syntaxique du C++ via listings -----------------------------
\usepackage{listings}

\definecolor{codebg}{rgb}{0.97,0.97,0.95}
\definecolor{codekw}{rgb}{0.10,0.10,0.60}
\definecolor{codecomment}{rgb}{0.20,0.50,0.20}
\definecolor{codestring}{rgb}{0.65,0.15,0.15}
\definecolor{coderule}{rgb}{0.80,0.80,0.80}

% Permet les caractères accentués (UTF-8) à l'intérieur des blocs de code.
\lstset{
    literate=%
        {é}{{\'e}}1 {è}{{\`e}}1 {ê}{{\^e}}1 {ë}{{\"e}}1
        {à}{{\`a}}1 {â}{{\^a}}1 {ä}{{\"a}}1
        {î}{{\^i}}1 {ï}{{\"i}}1
        {ô}{{\^o}}1 {ö}{{\"o}}1
        {ù}{{\`u}}1 {û}{{\^u}}1 {ü}{{\"u}}1
        {ç}{{\c c}}1 {É}{{\'E}}1 {È}{{\`E}}1 {À}{{\`A}}1
}

\lstdefinestyle{cpp}{
    language=C++,
    backgroundcolor=\color{codebg},
    basicstyle=\ttfamily\small,
    keywordstyle=\color{codekw}\bfseries,
    commentstyle=\color{codecomment}\itshape,
    stringstyle=\color{codestring},
    showstringspaces=false,
    numbers=left,
    numberstyle=\tiny\color{gray},
    numbersep=8pt,
    frame=single,
    rulecolor=\color{coderule},
    breaklines=true,
    tabsize=4,
    morekeywords={constexpr,noexcept,override,final,nullptr,auto,
        std,string,vector,size_t,uint64_t,int64_t,optional,namespace}
}
\lstset{style=cpp}

% Raccourci pour du code en ligne, en police machine à écrire.
\newcommand{\code}[1]{\texttt{#1}}

% Boîte « À retenir »
\usepackage{tcolorbox}
\newtcolorbox{aretenir}[1][]{
    colback=yellow!8, colframe=orange!70!black,
    title=#1, fonttitle=\bfseries, sharp corners, boxrule=0.5pt
}
\newtcolorbox{piege}[1][Piège]{
    colback=red!5, colframe=red!60!black,
    title=#1, fonttitle=\bfseries, sharp corners, boxrule=0.5pt
}


\title{\textbf{Chapitre 4 — Choisir et compter : optional, map, algorithmes}\\[0.3em]
       \large Modéliser l'absence, compter, et le projet de synthèse}
\author{Formation C++ moderne}
\date{}

\begin{document}
\maketitle

\begin{abstract}
\noindent Ce chapitre te donne trois briques qui reviennent partout dans le
code réel : quelques \textbf{algorithmes} de la STL (\code{<algorithm>},
\code{<numeric>}) pour trier, sommer et trouver des extrema ; \code{std::optional}
pour modéliser proprement « une valeur \emph{ou} rien » sans valeur-sentinelle
bricolée ; et \code{std::map} pour associer des clés à des valeurs et compter
sans effort. On les emploie ici \emph{en boîte noire} — juste ce qu'il faut pour
les utiliser correctement — et on les approfondira dans les chapitres dédiés.
À la fin, un \textbf{projet de synthèse} (un mini-REPL de statistiques) réunit
tout ce que tu as vu depuis le chapitre~1.
\end{abstract}

\tableofcontents
\bigskip\hrule\bigskip

% ===========================================================================
\section{\code{<algorithm>} et \code{<numeric>} : le strict nécessaire}

La bibliothèque standard fournit des dizaines d'algorithmes prêts à l'emploi.
On n'en prend ici que quatre, ceux dont tu auras besoin tout de suite. Ils
travaillent presque tous sur un \emph{intervalle} décrit par deux itérateurs :
\code{v.begin()} (premier élément) et \code{v.end()} (juste après le dernier).

\begin{aretenir}[En boîte noire pour l'instant]
On utilise ces algorithmes comme des outils tout faits, sans regarder dedans.
Le chapitre \emph{Algorithmes} les reprendra en profondeur (lambdas,
prédicats, complexité, ranges de C++20). Ici, objectif : savoir les
\emph{appeler} correctement.
\end{aretenir}

\subsection*{Trier : \code{std::sort}}

\code{std::sort(debut, fin)} réordonne les éléments de l'intervalle
\emph{sur place}, par ordre croissant par défaut.

\begin{lstlisting}
#include <algorithm>
#include <vector>

std::vector<int> v = {3, 1, 4, 1, 5, 9, 2, 6};
std::sort(v.begin(), v.end());     // v devient {1, 1, 2, 3, 4, 5, 6, 9}
\end{lstlisting}

Pour trier en ordre \emph{décroissant}, on passe un troisième argument : un
objet-comparateur. Le plus simple est \code{std::greater<>{}} (de
\code{<functional>}), qui compare avec « plus grand que » :

\begin{lstlisting}
#include <functional>   // pour std::greater

std::sort(v.begin(), v.end(), std::greater<>{});   // {9, 6, 5, 4, 3, 2, 1, 1}
\end{lstlisting}

\begin{piege}[\code{std::sort} modifie le conteneur]
Le tri se fait \textbf{en place} : l'ordre d'origine est perdu. Si tu as besoin
de garder la version non triée, fais une copie d'abord
(\code{auto copie = v;}) et trie la copie.
\end{piege}

\subsection*{Sommer : \code{std::accumulate}}

\code{std::accumulate(debut, fin, valeur\_initiale)} (de \code{<numeric>}) part
de la valeur initiale et y ajoute chaque élément. Il renvoie la somme.

\begin{lstlisting}
#include <numeric>

std::vector<int> notes = {12, 15, 9, 18};
int total = std::accumulate(notes.begin(), notes.end(), 0);   // 54
\end{lstlisting}

\begin{piege}[Le type de l'accumulateur décide du calcul]
Le type du \emph{troisième argument} fixe le type interne de l'addition. Avec
\code{0} (un \code{int}), tout est calculé en \code{int}, même sur un vecteur de
\code{double} : tu perds les décimales !

\begin{lstlisting}
std::vector<double> v = {1.5, 2.5, 3.0};
double faux = std::accumulate(v.begin(), v.end(), 0);     // 6.0 ? NON : 6 (tronqué)
double bon  = std::accumulate(v.begin(), v.end(), 0.0);   // 7.0, correct
\end{lstlisting}

\noindent Règle : pour sommer des \code{double}, démarre à \code{0.0}. Pour des
entiers qui pourraient déborder, démarre à \code{0LL} (un \code{long long}).
\end{piege}

\subsection*{Extrema : \code{std::min\_element} et \code{std::max\_element}}

Ces deux-là renvoient un \textbf{itérateur} vers le plus petit (resp. plus
grand) élément — \emph{pas} la valeur elle-même. Pour obtenir la valeur, on
déréférence l'itérateur avec \code{*}.

\begin{lstlisting}
#include <algorithm>

std::vector<int> v = {3, 1, 4, 1, 5, 9, 2};
auto it_min = std::min_element(v.begin(), v.end());   // itérateur, pas un int
int mini    = *it_min;                                 // 1  (on déréférence)
int maxi    = *std::max_element(v.begin(), v.end());   // 9  (directement)
\end{lstlisting}

\begin{piege}[Conteneur vide $\to$ itérateur invalide]
Sur un intervalle vide, \code{min\_element}/\code{max\_element} renvoient
\code{end()}, qui ne pointe sur rien : le déréférencer est un \emph{comportement
indéfini}. Vérifie \code{if (!v.empty())} avant. (C'est précisément un cas où
\code{std::optional} de la section suivante sera utile.)
\end{piege}

\begin{aretenir}[Les quatre en une phrase]
\code{std::sort(b, e[, comp])} trie en place ;
\code{std::accumulate(b, e, init)} somme (soigne le type de \code{init}) ;
\code{std::min\_element}/\code{max\_element} renvoient un \emph{itérateur} à
déréférencer. Tous travaillent sur \code{begin()}/\code{end()}.
\end{aretenir}

% ===========================================================================
\section{\code{std::optional<T>} : une valeur ou rien}

Comment une fonction dit-elle « je n'ai pas de résultat » ? En venant d'autres
langages, le réflexe est une \emph{valeur-sentinelle} : renvoyer \code{-1}, ou
une chaîne vide, ou \code{NaN}. C'est fragile : il faut se souvenir de la
convention, et parfois la sentinelle est une valeur légitime (que faire si
\code{-1} est un résultat possible ?).

\code{std::optional<T>} (de \code{<optional>}) résout ça proprement : il
contient \emph{soit} une valeur de type \code{T}, \emph{soit} rien du tout
(\code{std::nullopt}). Le « rien » est explicite, distinct de toute valeur réelle.

\subsection*{Produire un \code{optional}}

\begin{lstlisting}
#include <optional>

std::optional<int> chercher(int x) {
    if (x < 0)
        return std::nullopt;   // "rien" : pas de résultat valable
    return x * 2;              // une valeur : conversion implicite en optional<int>
}
\end{lstlisting}

Tu renvoies \code{std::nullopt} pour l'absence, et directement une valeur de
type \code{T} sinon — elle est automatiquement « emballée » dans l'optional.

\subsection*{Consommer un \code{optional}}

Trois façons de regarder ce qu'il contient :

\begin{lstlisting}
std::optional<int> r = chercher(21);

if (r.has_value()) {           // y a-t-il une valeur ?
    int v = r.value();         // .value() la récupère (lève si vide)
}

if (r) {                       // idem, plus court : un optional se teste comme un bool
    int v = *r;                // *r déréférence, comme un pointeur
}

int v = r.value_or(0);         // la valeur, ou 0 si l'optional est vide
\end{lstlisting}

\begin{piege}[\code{.value()} sur un optional vide \emph{lève}]
Appeler \code{.value()} (ou déréférencer \code{*r}) alors que l'optional est
vide lance une exception \code{std::bad\_optional\_access} (pour \code{.value()})
ou est un comportement indéfini (pour \code{*r}). \textbf{Teste toujours} avant :
\code{if (r) \{ ... \}}, ou utilise \code{r.value\_or(defaut)}.
\end{piege}

\subsection*{Exemple complet : \code{safe\_divide}}

La division par zéro n'a pas de résultat. Plutôt qu'un \code{-1} ambigu ou une
exception, on renvoie « rien » :

\begin{lstlisting}
#include <optional>

std::optional<double> safe_divide(double a, double b) {
    if (b == 0.0)
        return std::nullopt;   // pas de résultat : division impossible
    return a / b;
}

// Utilisation :
if (auto q = safe_divide(10.0, 4.0)) {   // init-statement + test en une ligne
    std::cout << "Quotient = " << *q << '\n';   // 2.5
} else {
    std::cout << "Division par zéro\n";
}
\end{lstlisting}

\begin{aretenir}[Pourquoi mieux qu'une sentinelle ?]
Avec \code{std::optional}, l'absence est dans le \emph{type} : impossible de
l'oublier, le compilateur t'oblige à extraire la valeur. Aucune valeur réelle
n'est « réservée » pour signaler l'erreur. Le code qui appelle est clair :
\code{if (resultat) \{ ... \}}.
\end{aretenir}

% ===========================================================================
\section{\code{std::map<K,V>} : un dictionnaire trié par clé}

\code{std::map<Cle, Valeur>} (de \code{<map>}) associe des \emph{clés} uniques à
des \emph{valeurs}. Si tu viens de Python, c'est le \code{dict} — avec une
différence majeure : la map est en permanence \textbf{triée par clé}.

\begin{lstlisting}
#include <map>
#include <string>

std::map<std::string, int> ages;
ages["Alice"] = 30;            // insère la paire ("Alice", 30)
ages["Bob"]   = 25;
int a = ages["Alice"];         // 30  (accès par clé)
std::size_t n = ages.size();   // 2   (nombre de paires)
\end{lstlisting}

\subsection*{Compter sans effort : \code{++m[cle]}}

C'est l'idiome vedette de la map. L'opérateur \code{[]} a une propriété très
utile : si la clé \emph{n'existe pas encore}, il l'insère avec une valeur
\textbf{par défaut} (\code{0} pour un \code{int}), puis renvoie une référence
vers elle. Donc :

\begin{lstlisting}
std::map<std::string, int> compte;
++compte["chat"];    // "chat" absent -> créé à 0, puis incrémenté -> 1
++compte["chat"];    // "chat" présent -> 1, puis incrémenté        -> 2
++compte["chien"];   // "chien" absent -> créé à 0, puis incrémenté -> 1
\end{lstlisting}

Pas besoin de tester « la clé existe-t-elle ? » : \code{++compte[mot]} fait tout.

\subsection*{Parcourir : structured bindings, et c'est trié}

On parcourt une map avec une boucle \code{for} par intervalle. Chaque élément
est une \emph{paire} (clé, valeur) ; les \textbf{structured bindings}
(C++17) la décomposent en deux noms d'un coup :

\begin{lstlisting}
for (const auto& [mot, n] : compte) {    // [cle, valeur] : décomposition
    std::cout << mot << " : " << n << '\n';
}
\end{lstlisting}

\begin{aretenir}[Le parcours d'une map donne les clés \emph{dans l'ordre}]
Contrairement à un \code{dict} Python (ordre d'insertion) ou à une table de
hachage, parcourir une \code{std::map} visite toujours les clés \textbf{triées}
(ordre croissant). Pour l'exemple ci-dessus, la sortie sera \code{chat} puis
\code{chien}, quel que soit l'ordre d'insertion. Pratique pour produire des
résultats déterministes et lisibles.
\end{aretenir}

\subsection*{Exemple complet : compter des mots}

\begin{lstlisting}
#include <iostream>
#include <map>
#include <string>

int main() {
    std::map<std::string, int> freq;
    std::string mot;
    while (std::cin >> mot) {     // lit mot à mot jusqu'à la fin de l'entrée
        ++freq[mot];             // compte chaque mot
    }
    for (const auto& [m, n] : freq) {   // affichage trié par mot
        std::cout << m << " : " << n << '\n';
    }
}
\end{lstlisting}

\begin{piege}[\code{m[cle]} en lecture \emph{insère} la clé]
Attention : \code{compte["xyz"]} insère \code{"xyz"} (avec la valeur 0) même si
tu voulais seulement \emph{lire}. Pour vérifier la présence sans insérer, utilise
\code{compte.contains("xyz")} (C++20) ou \code{compte.find("xyz") != compte.end()}.
\end{piege}

% ===========================================================================
\section{Le modèle « spirale »}

\begin{aretenir}[On revient sur tout en plus profond]
Ce chapitre suit un modèle \textbf{en spirale} : on découvre \code{std::map},
\code{std::optional} et les algorithmes \emph{en boîte noire}, avec juste assez
de détails pour les employer sans se tromper. On ne plonge pas encore dans leur
mécanique interne (arbre équilibré sous la map, gestion mémoire de l'optional,
itérateurs et lambdas des algorithmes). Les chapitres ultérieurs —
\emph{Conteneurs} et \emph{Algorithmes} — repasseront par ces mêmes outils pour
en détailler le fonctionnement, le coût et les variantes. Tu les utilises
aujourd'hui ; tu les comprendras de fond en comble bientôt.
\end{aretenir}

% ===========================================================================
\section{Récapitulatif}

\begin{itemize}[noitemsep]
    \item \code{std::sort(b, e)} trie en place ; ajoute \code{std::greater<>{}}
          pour l'ordre décroissant.
    \item \code{std::accumulate(b, e, init)} somme — le type de \code{init}
          (\code{0.0}, \code{0LL}\ldots) gouverne le calcul.
    \item \code{std::min\_element}/\code{max\_element} renvoient un
          \emph{itérateur} ; déréférence avec \code{*}.
    \item \code{std::optional<T>} : \code{std::nullopt} pour rien, une valeur
          sinon ; teste avec \code{if (opt)} ; lis avec \code{*opt} /
          \code{opt.value()} / \code{opt.value\_or(...)}.
    \item \code{std::map<K,V>} : \code{++m[cle]} pour compter ; parcours
          \code{for (const auto\& [k, v] : m)} \textbf{trié par clé}.
\end{itemize}

\bigskip\hrule\bigskip
\noindent\textbf{À toi de jouer.} Les exercices du dossier \code{exercices/}
mettent ces briques en pratique, de la plus simple à la plus complète :

\begin{itemize}[noitemsep]
    \item \textbf{optionnel} — une fonction qui renvoie \code{std::optional<T>}
          (par exemple une recherche, un parsing qui peut échouer).
    \item \textbf{compteur-mots} (autonome) — un programme complet,
          \code{main} compris, qui compte les occurrences via une
          \code{std::map} et affiche le résultat trié.
    \item \textbf{moyenne-ou-rien} — calculer la moyenne d'un \code{vector},
          mais renvoyer \code{std::nullopt} si le vecteur est vide (combine
          \code{accumulate} et \code{optional}).
    \item \textbf{config-parser} — lire des lignes \code{cle = valeur} et les
          ranger dans une \code{std::map}, en renvoyant un \code{optional} quand
          on interroge une clé absente.
    \item \textbf{PROJET de synthèse : mini-REPL de statistiques} — le gros
          morceau. Une boucle de lecture-évaluation-affichage qui accepte des
          commandes (ajouter des nombres, demander somme, moyenne, min, max,
          tri\ldots), réunissant \emph{tout} depuis le chapitre~1 :
          \code{std::vector}, \code{std::string} et parsing de flux, les
          algorithmes de ce chapitre, \code{std::optional} pour les résultats
          indéfinis, et \code{std::map} pour, par exemple, un historique de
          commandes.
\end{itemize}

\begin{aretenir}[Le projet sera « cmakifié » au chapitre suivant]
Tu écriras le mini-REPL en compilation manuelle, dans l'esprit des chapitres
précédents. Au \textbf{chapitre suivant (CMake)}, on reprendra ce même projet
pour le construire proprement avec CMake et le brancher sur les tests
automatisés — première étape vers un vrai flux de build de projet.
\end{aretenir}

\noindent Bon courage !

\end{document}
